% CrashVLM: workshop-style draft. Neutral two-column layout; port to the venue template
% (CVPR/ICCV workshop, NeurIPS workshop) before submission.
% Numbers come from paper/numbers.tex and paper/tables/*.tex, generated by
% scripts/paper_tables.py. Do not type results into this file.
\documentclass[10pt,twocolumn]{article}
\usepackage[margin=0.75in]{geometry}
\usepackage{booktabs,amsmath,graphicx,xcolor,hyperref}
\usepackage[numbers,sort&compress]{natbib}
\newcommand{\pending}{\textcolor{orange}{\textit{pend.}}}
\newcommand{\todo}[1]{\textcolor{red}{[#1]}}
% generated by scripts/paper_tables.py
\newcommand{\NCatalog}{1583}
\newcommand{\NTests}{100}
\newcommand{\NClips}{200}
\newcommand{\NTestTests}{23}
\newcommand{\NVerified}{193}
\newcommand{\NLabelled}{200}
\newcommand{\BlindAirbagKnown}{1.00}
\newcommand{\BlindHeadKnown}{0.94}
\newcommand{\BlindAirbagUnknown}{0.75}
\newcommand{\ProbeAirbagUnknown}{0.60}
\newcommand{\BlindImpactUnknown}{0.75}
\newcommand{\ResnetImpactUnknown}{0.75}
\newcommand{\LargeProbeAirbagUnknown}{0.70}
\newcommand{\VlmAirbagUnknown}{0.75}
\newcommand{\VlmHeadUnknown}{0.39}
\newcommand{\VlmImpactUnknown}{0.20}
\newcommand{\VlmParse}{1.00}
\newcommand{\VlmTestType}{0.37}
\newcommand{\VlmContradict}{0.11}
\newcommand{\VlmLatency}{104}
\newcommand{\BlindImpactEight}{0.35}
\newcommand{\ZsImpactEight}{0.20}
\newcommand{\ZsImpactEightDiff}{-0.15}
\newcommand{\VlImpactEight}{0.20}
\newcommand{\VlImpactEightDiff}{-0.15}
\newcommand{\LoraImpactEight}{0.65}
\newcommand{\LoraImpactEightDiff}{+0.30}
\newcommand{\BlindImpactSixteen}{0.75}
\newcommand{\ZsImpactSixteen}{0.20}
\newcommand{\ZsImpactSixteenDiff}{-0.55}
\newcommand{\VlImpactSixteen}{0.40}
\newcommand{\VlImpactSixteenDiff}{-0.35}
\newcommand{\LoraImpactSixteen}{0.90}
\newcommand{\LoraImpactSixteenDiff}{+0.15}
\newcommand{\BlindAirbagEight}{0.25}
\newcommand{\ZsAirbagEight}{0.70}
\newcommand{\ZsAirbagEightDiff}{+0.45}
\newcommand{\VlAirbagEight}{0.90}
\newcommand{\VlAirbagEightDiff}{+0.65}
\newcommand{\LoraAirbagEight}{0.80}
\newcommand{\LoraAirbagEightDiff}{+0.55}
\newcommand{\BlindAirbagSixteen}{0.75}
\newcommand{\ZsAirbagSixteen}{0.75}
\newcommand{\ZsAirbagSixteenDiff}{+0.00}
\newcommand{\VlAirbagSixteen}{0.90}
\newcommand{\VlAirbagSixteenDiff}{+0.15}
\newcommand{\LoraAirbagSixteen}{0.90}
\newcommand{\LoraAirbagSixteenDiff}{+0.15}
\newcommand{\BlindHeadEight}{0.22}
\newcommand{\ZsHeadEight}{0.33}
\newcommand{\ZsHeadEightDiff}{+0.11}
\newcommand{\VlHeadEight}{0.11}
\newcommand{\VlHeadEightDiff}{-0.11}
\newcommand{\LoraHeadEight}{0.56}
\newcommand{\LoraHeadEightDiff}{+0.33}
\newcommand{\BlindHeadSixteen}{0.67}
\newcommand{\ZsHeadSixteen}{0.39}
\newcommand{\ZsHeadSixteenDiff}{-0.28}
\newcommand{\VlHeadSixteen}{0.17}
\newcommand{\VlHeadSixteenDiff}{-0.50}
\newcommand{\LoraHeadSixteen}{0.83}
\newcommand{\LoraHeadSixteenDiff}{+0.17}
\newcommand{\LoraSteps}{20}
\newcommand{\LoraPlannedSteps}{183}
\newcommand{\LoraTrainHours}{6.5}
\newcommand{\VlTestType}{0.48}
\newcommand{\VlContradict}{0.25}
\newcommand{\VlFire}{0.07}
\newcommand{\LoraTestType}{0.46}
\newcommand{\LoraContradict}{0.09}
\newcommand{\LoraFire}{0.00}
\newcommand{\NPhysics}{94}
\newcommand{\NFlowHeldOut}{12}
\newcommand{\FlowHeldOutWithinEight}{12}
\newcommand{\FlowTestMedErr}{0.9}
\newcommand{\FlowTrainRho}{0.71}
\newcommand{\BlindTrainRho}{-0.12}
\newcommand{\FlowBias}{+3.1}


\title{CrashVLM: How Much Do Vision-Language Models See in\\Crash Test Videos, and How Much Is the Clock?}
\author{Muhammad Faheem Arshad}
\date{Draft, \today}

\begin{document}
\maketitle

\begin{abstract}
Crash test videos are recorded for every vehicle safety rating, but they are evaluated by hand.
We study whether open vision-language models (VLMs) can detect and time crash-relevant events,
such as impact, airbag deployment, head contact and glass failure, in high-speed NCAP footage
from the public NHTSA database. We release a reproducible pipeline that selects
\NTests{} tests (\NClips{} camera views) from \NCatalog{} candidates. It removes burned-in
overlays that would otherwise leak the answer and recovers each clip's physical time axis. Our
main finding is methodological. Because NCAP protocols are highly standardised, a blind predictor
that never sees the video times airbag deployment within 16\,ms in \BlindAirbagKnown{} of test
cases once the laboratory clock is known. Timing benchmarks on standardised crash tests are
therefore only informative without the clock. In that setting, linear probes on frozen SigLIP~2
features beat the blind baseline (airbag: \ProbeAirbagUnknown{} vs.\ \BlindAirbagUnknown{} hit@16),
while zero-shot text prompting times events poorly. Grounding the video in the vehicles'
own accelerometers, we derive the measured time of maximum crush for \NPhysics{} tests and show
that clock-free optical flow on the vehicle body recovers it within 8\,ms in
\FlowHeldOutWithinEight{} of \NFlowHeldOut{} held-out frontal tests, where visual labels and
appearance probes fail. We further report LoRA-adapted video VLMs
producing structured crash reports \pending{}, a hallucination analysis, and a label-quality
protocol for model-produced annotations.
\end{abstract}

\section{Introduction}
Vehicle crash tests produce two kinds of evidence: instrumented measurements (accelerations,
injury criteria, intrusions) and high-speed video from a dozen cameras. The measurements feed
the ratings. The video is inspected by engineers to understand \emph{what happened when}: did
the airbag deploy before head contact, did the door open, which glass failed. Automating this
reading would make large test archives searchable and, eventually, link physical tests to
simulation. Two developments make this timely: open VLMs now accept long image
sequences~\citep{bai2025qwen3vl,wang2025internvl35}, and CLIP-type
encoders~\citep{radford2021clip,zhai2023siglip,tschannen2025siglip2} give strong frozen frame
features.

We ask: \emph{Can open VLMs, adapted with a small annotated dataset, reliably detect and
describe crash events in crash test video?} Answering it honestly required three things that
turned out to matter more than the choice of model:
(i)~removing burned-in overlays (timestamps, frame counters, a red$\to$green impact lamp) that
leak the labels;
(ii)~recovering each clip's time axis from the overlay itself, because frame rate and start
time vary between laboratories; and
(iii)~comparing every model against a \emph{blind} baseline that exploits test standardisation,
a shortcut in the sense of~\citet{geirhos2020shortcut}.

\paragraph{Contributions.}
\begin{itemize}\setlength\itemsep{0pt}
\item An open, scripted pipeline over the public NHTSA database: catalogue, vehicle-grouped
split, overlay removal with an automated leak check, and OCR-based time recovery.
\item A frame-level event ontology and \NLabelled{} annotated views, with a documented
model-annotator protocol and cross-model agreement checks.
\item Evidence that timing benchmarks on standardised crash tests are dominated by the clock,
and a \emph{t0-unknown} protocol that removes it.
\item Baselines from blind to probe to zero-shot VLM to LoRA-adapted VLM, with hallucination and
efficiency analysis \pending{}.
\item Physics ground truth from NHTSA sensor data and a clock-free video-kinematics method
validated against it.
\end{itemize}

\section{Related work}
\textbf{Image-text encoders} such as CLIP~\citep{radford2021clip} and SigLIP~\citep{zhai2023siglip,tschannen2025siglip2}
support zero-shot classification by text prompts; we use them as frozen frame encoders.
\textbf{Video VLMs.} Qwen2.5-VL~\citep{bai2025qwen25vl}, Qwen3-VL~\citep{bai2025qwen3vl} and
InternVL3.5~\citep{wang2025internvl35} process frame sequences. Time-sensitive variants
target temporal grounding~\citep{ren2023timechat}, a task framed as temporal activity
localisation by language query in~\citet{gao2017tall}. Parameter-efficient adaptation (LoRA,
QLoRA~\citep{hu2021lora,dettmers2023qlora}) makes fine-tuning feasible on a single consumer GPU.
Object hallucination in VLMs is measured by probing for absent objects~\citep{li2023pope}; our
fire/leak probe follows the same idea.
\textbf{Machine learning for crashworthiness} has concentrated on simulation surrogates:
graph networks on meshes~\citep{pfaff2020meshgraphnets,kneifl2024multihierarchical,li2025gnncrash},
accelerated crash-dynamics models~\citep{nabian2025crashdynamics}, random-feature
surrogates~\citep{herold2022randomfeatures} and large simulation datasets~\citep{elrefaie2026carcrashnet}.
An arXiv search did not find prior work that applies VLMs to physical crash test video. This
work is a step towards connecting the two.
\textbf{Classic video models} such as R(2+1)D~\citep{tran2017r21d} and image CNNs~\citep{he2015resnet}
serve as non-VLM references.

\section{Data}
\textbf{Source and licence.} The NHTSA Vehicle Crash Test Database is publicly accessible
without registration. data.gov lists its licence as ``unknown'', and the footage is produced by
contract laboratories, so we redistribute only code, test identifiers, annotations and derived
metrics. Everything else is re-downloaded by the scripts.
\textbf{Selection.} From the \NCatalog{} NCAP frontal-barrier, side-barrier and side-pole tests
(2012 onwards) we draw a seeded sample of \NTests{} tests. The split is by vehicle make and model
(no vehicle in two splits) and balanced per test type, giving \NTestTests{} test-split tests.
Each test contributes an \emph{exterior} and an \emph{occupant} view.
\textbf{Overlay removal.} Laboratories burn in time, frame counters and a red lamp that turns
green at impact. We crop rendered border panels, where a line is kept as panel only if it is
uniform \emph{and} at least 90\,\% static over the clip; this rejects dark cabin areas that
uniformity alone would crop. We then mask OCR text boxes that behave like overlays and lamp-like
blobs that switch from red to green once. A final check scans every processed clip for surviving
lamps.
\textbf{Time axis.} Numbers read from the overlay are fitted as $t = a\,\mathrm{frame}+b$ with a
robust Theil--Sen fit. Clips without a printed time use the printed impact frame (T0) and frame
rate. \NVerified{} of \NClips{} clips have a self-consistent fit; only these contribute timing
metrics.

\section{Annotation and label quality}
Seven point events are tagged on the first visible frame at 8\,ms resolution (impact, airbag
deployment, maximum deformation, head contact, windshield crack, side-glass break, door
opening), together with a per-clip summary. \textbf{All \NLabelled{} views were annotated by a
language model (Claude), not by a human}, from time-stamped contact sheets under a written
guideline; the protocol is disclosed as a limitation. Quality is estimated by
(a)~two blind checks by a different model of the same family (Claude Sonnet), 75 views, and
(b)~a blind repeat by an independent instance of the primary model. Timed events agree closely,
whereas glass-type and structural attributes do not (details and all numbers in the repository,
\texttt{docs/ANNOTATION\_QA.md}). Labels are therefore used in three reliability tiers, and
only tier~A (impact, airbag, head contact, glass events) enters the headline metrics.
Same-model agreement is much higher than cross-model agreement on subjective labels. We treat
it as a measure of reproducibility, not correctness.

\section{Methods}
\textbf{Blind baseline.} The median training onset per (event, test type, view); the event is
predicted absent if it occurs in fewer than half of the training clips.
\textbf{Probes.} Per event, a logistic regression on frozen frame embeddings predicts
``has the event happened yet''. The onset is the start of the first run of at least two frames
above a threshold; C and the threshold are chosen on the validation split.
\textbf{Zero-shot.} SigLIP~2 scores each frame against positive and negative prompts and the
same decoder is applied.
\textbf{VLM reports.} Qwen3.5-4B (bf16) and Qwen3-VL-8B (4-bit) receive 24 frames at 8\,ms
steps, each labelled with its time since the \emph{first frame shown}, and must return a JSON
report: test type, event times and summary fields. A strict parser records every deviation.
\textbf{LoRA.} Rank-16 adapters on the language model's attention are trained on
instruction examples built from the labels, with four random start crops per training clip and
the loss on answer tokens only.

\section{Evaluation protocol}
Metrics per event are presence F1 and hit@$k$: the share of timing-eligible positives predicted
within $k$\,ms, where a missed detection counts as a miss. 95\,\% intervals come from resampling
\emph{tests}. We report two settings.
In \textbf{t0 known}, each clip's time axis is anchored at the laboratory impact time.
In \textbf{t0 unknown}, each test clip is cropped at a seeded random start (0--40\,ms, per clip
identifier) and predictions may use only time since the first remaining frame, as for any video
in the wild. Video models are unaffected, since they read pixels, not the clock. The blind
baseline must guess where in the clip the event falls.

\section{Results}
\begin{table}[t]\centering\small
\caption{Hit@16 on the test split, \textbf{t0 known}. The blind baseline uses no video.}
\label{tab:known}
\resizebox{\columnwidth}{!}{\begin{tabular}{lcccc}
\toprule
Method & Airbag & Head contact & Windshield crack & Side glass \\
\midrule
Blind (no video) & 1.00 & 0.94 & 1.00 & 0.36 \\
ResNet-50 + probe & 0.40 & 0.44 & 0.50 & 0.27 \\
SigLIP 2-B zero-shot & 0.40 & 0.28 & 0.25 & 0.00 \\
SigLIP 2-B + probe & 0.55 & 0.50 & 0.62 & 0.27 \\
SigLIP 2-So400m zero-shot & \pending & \pending & \pending & \pending \\
SigLIP 2-So400m + probe & \pending & \pending & \pending & \pending \\
\bottomrule
\end{tabular}
}
\end{table}
\begin{table}[t]\centering\small
\caption{Hit@16 on the test split, \textbf{t0 unknown} (random start crops). ``--'': no prompt
defined for that event.}
\label{tab:unknown}
\resizebox{\columnwidth}{!}{\begin{tabular}{lccccc}
\toprule
Method & Impact & Airbag & Head contact & Windshield crack & Side glass \\
\midrule
Blind (no video) & 0.45 & 0.35 & 0.33 & 0.50 & 0.09 \\
ResNet-50 + probe & 0.75 & 0.45 & 0.44 & 0.50 & 0.27 \\
SigLIP 2-B zero-shot & -- & 0.45 & 0.33 & 0.25 & 0.00 \\
SigLIP 2-B + probe & 0.55 & 0.60 & 0.50 & 0.62 & 0.27 \\
SigLIP 2-So400m zero-shot & -- & 0.35 & 0.28 & 0.12 & 0.00 \\
SigLIP 2-So400m + probe & 0.45 & 0.70 & 0.50 & 0.62 & 0.45 \\
\bottomrule
\end{tabular}
}
\end{table}

\textbf{The clock gives the answer away (Tab.~\ref{tab:known}).} With the laboratory clock,
the blind baseline reaches \BlindAirbagKnown{} (airbag) and \BlindHeadKnown{} (head contact)
hit@16 and beats every video model. NCAP fixes impact speed, barrier and dummy, so event times
barely vary between vehicles at 8\,ms resolution. Reported timing accuracy on such data is
uninterpretable without this baseline.

\textbf{Without the clock, video helps (Tab.~\ref{tab:unknown}).} Linear probes beat the blind
baseline for every event, most clearly for finding the impact itself (ResNet-50:
\ResnetImpactUnknown{} vs.\ \BlindImpactUnknown{}). Most intervals overlap at
\NTestTests{} test tests. Zero-shot prompting detects airbags but times them poorly and fails on
glass. The larger SigLIP~2 encoder is \pending{}.

\textbf{More labels barely help frozen probes (Tab.~\ref{tab:lc}).} From 25\,\% to 100\,\% of
the training tests, hit@16 stays within seed noise except for side-glass breakage. The frozen
features and the onset decoder, not label volume, limit the probes.

\begin{table}[t]\centering\small
\caption{Learning curve of the SigLIP 2-B probes (t0 unknown, hit@16, mean $\pm$ std over
three subsets).}\label{tab:lc}
\begin{tabular}{lccc}
\toprule
Event & 25\% & 50\% & 100\% \\
\midrule
Airbag & 0.53 $\pm$ 0.08 & 0.43 $\pm$ 0.03 & 0.60 \\
Head contact & 0.44 $\pm$ 0.11 & 0.48 $\pm$ 0.18 & 0.50 \\
Impact & 0.42 $\pm$ 0.03 & 0.55 $\pm$ 0.05 & 0.55 \\
Side glass & 0.03 $\pm$ 0.05 & 0.18 $\pm$ 0.00 & 0.27 \\
Windshield crack & 0.50 $\pm$ 0.00 & 0.50 $\pm$ 0.12 & 0.62 \\
\bottomrule
\end{tabular}

\end{table}

\textbf{Structured VLM reports (Tab.~\ref{tab:vlm}).} \pending{} Zero-shot and LoRA-adapted
Qwen reports, scored on the same test clips and crops: event timing (comparable to
Tab.~\ref{tab:unknown}), test-type accuracy against NHTSA metadata, and hallucination rates.

\begin{table}[t]\centering\small
\caption{VLM report quality on test clips. Fire hallucination: share of reports claiming fire or
a leak (none occurs in any clip). Contradicting claims: reported events that the labels mark as
absent.}\label{tab:vlm}
\resizebox{\columnwidth}{!}{\begin{tabular}{lccccc}
\toprule
Model & Parse & Test type & Fire halluc. & Contradicting claims & Latency (s) \\
\midrule
qwen3.5-4b\_zs & \pending & \pending & \pending & \pending & \pending \\
qwen3-vl-8b-4bit\_zs & \pending & \pending & \pending & \pending & \pending \\
qwen3.5-4b\_lora & \pending & \pending & \pending & \pending & \pending \\
\bottomrule
\end{tabular}
}
\end{table}

\textbf{Severity.} The annotated visual severity does not correlate with measured crush or driver
HIC beyond chance (repository, \texttt{reports/m6/severity.md}). Whether VLM severity ratings
track the measurements is \pending{}.

\section{Physics grounding}
NHTSA publishes the instrumented channels of each test (ISO-MME). We integrate structural
accelerations from the lab-measured impact speed. For frontal and pole tests, maximum dynamic
crush is where the velocity towards the obstacle reaches zero. For side tests with the moving
barrier, it is where barrier and vehicle velocities meet; for a yawing vehicle that never quite
meets the barrier velocity, we take the minimum of the closing velocity. This gives a
\emph{measured} time of maximum crush for \NPhysics{} of \NTests{} tests. Dynamic crush exceeds
NHTSA's residual crush, as rebound implies (\texttt{reports/physics/checks.md}).

\textbf{Visual labels vs.\ physics.} Against the measured times, the annotated
``maximum deformation'' frames are no closer than a constant per test type, which confirms that
this label cannot be read reliably from the video by eye.

\textbf{Video kinematics.} Frame-appearance probes trained on the measured time fail (below the
blind baseline). Maximum crush is a kinematic turning point, not an appearance. Tracking the
body instead does work. We run pyramidal Lucas--Kanade optical flow on the rear part of the
visible body in the frontal side view and find where its horizontal velocity stops. The method
is clock-free, with one bias of \FlowBias{}\,ms estimated on training tests. Designed on the
training split, it tracks the vehicle-to-vehicle variation of the measured time (Spearman
\FlowTrainRho{} vs.\ \BlindTrainRho{} for the blind constant). On the \NFlowHeldOut{} held-out
frontal tests it is within 8\,ms of the accelerometer time in \FlowHeldOutWithinEight{} cases,
with a median test error of \FlowTestMedErr{}\,ms, including an early outlier the constant
misses. It is the one vision result in this work validated against a physical measurement
rather than a label. The scope is small (frontal view, 12 tests).

\section{Limitations}
\begin{itemize}\setlength\itemsep{0pt}
\item \textbf{Model-produced labels.} All annotations come from a language model. Cross-model
checks reduce but cannot exclude shared errors, and a human-verified subset is the most
important missing piece.
\item \textbf{Small test set.} \NTestTests{} tests: most differences between video methods are
within the confidence intervals.
\item \textbf{Narrow domain.} US NCAP only: fixed protocols, five laboratories, and camera setups
that can identify the laboratory.
\item \textbf{8\,ms resolution} and annotator clustering on tiles; no claims below 8\,ms.
\item \textbf{Physics grounding is partial.} Sensor targets exist for one event (maximum crush); video kinematics are validated for frontal views only.
\item \textbf{No geometry yet.} The simulation and geometric-deep-learning link
(PointNet++~\citep{qi2017pointnetpp}, mesh GNNs~\citep{pfaff2020meshgraphnets}) is future work.
\end{itemize}

\section{Conclusion and outlook}
Reading crash test video with open models is feasible, but benchmarks on standardised tests must
remove the clock before model comparisons mean anything. The next steps are:
\begin{itemize}\setlength\itemsep{0pt}
\item extend video kinematics to side and pole views and to crush \emph{depth} (scaled flow vs.\ integrated acceleration), via the photogrammetry targets on NCAP vehicles;
\item add a human-verified label subset;
\item connect video evidence to simulation, by rendering explicit-solver runs of parametric
crash structures and asking whether the same models transfer from simulation to physical
video.
\end{itemize}

\paragraph{Candidate venues.} A computer-vision workshop on autonomous driving or industrial
vision at CVPR, ICCV or ECCV (check the current year's list); the NeurIPS Datasets \&
Benchmarks track (for the protocol and data pipeline); IEEE IV or ITSC; and crash-safety venues:
the ESV Conference, SAE WCX, and the Stapp Car Crash Conference (after adding a human-verified
subset).

{\small\bibliographystyle{plainnat}\bibliography{refs}}
\end{document}
